\documentclass[a4paper, 12pt]{article}

\usepackage[utf8]{inputenc}
\usepackage{enumitem}
\usepackage{amsmath}
\usepackage{amssymb}

\title{E.B.1.3}
\author{Lorenzo Burello}
\date{\today}

\begin{document}
\maketitle

\section{Definizione del problema}
\subsection{Consegna}
Modellare, mediante una formula $\phi$ in logica del primo ordine,
la seguente affermazione, scegliendo opportuni insiemi di simboli
di predicato e/o funzione: \newline
\begin{quote}
    I cani da caccia abbaiano di notte. Chi possiede gatti non può avere topi in casa.
    Chi ha il sonno leggero non può possedere nessun animale che abbaia di notte.
    Aldo possiede un gatto oppure un cane da caccia.
\end{quote}
\subsection{Definizione dell'alfabeto}
$\mathcal{P}$:
\begin{itemize}[label=-]
    \item CaneDaCaccia/1
    \item Gatto/1
    \item Persona/1
    \item HaTopiInCasa/1
    \item HaIlSonnoLeggero/1
    \item AbbaiaDiNotte/1
    \item Possiede/2
\end{itemize}

$\mathcal{F}$:
\begin{itemize}[label=-]
    \item Aldo/0
\end{itemize}

\section{Creazione di $\phi$}
\subsection{Creazione di sotto-formule}
\begin{quote}
    \textbf{I cani da caccia abbaiano di notte}
\end{quote}
\[\phi_1 := \forall c \;\; CaneDaCaccia(c) \rightarrow AbbaiaDiNotte(c)\]
\begin{quote}
    \textbf{chi possiede gatti non può avere topi in casa}
\end{quote}
\[\phi_2 := \forall x,g\;\; (Possiede(x, g) \land Gatto(g)) \rightarrow \lnot HaTopiInCasa(x) \]
\begin{quote}
    \textbf{Chi ha il sonno leggero non può possedere nessun animale che abbaia di notte}
\end{quote}
\[ \phi_3 := \forall x \;\; (Persona(x) \land HaIlSonnoLeggero(x))\rightarrow \nexists a \;\; Animale(a) \land \]
\[\land AbbaiaDiNotte(a) \land (Possiede(x, a))\]
\begin{quote}
    \textbf{Aldo possiede un gatto o un cane da caccia}
\end{quote}
\[ \phi_4 := \exists a\;\; Possiede(Aldo, a) \land (Gatto(a) \lor CaneDaCaccia(a))\]
\subsection{Unione delle formule}
\[ \phi = \phi_1 \land \phi_2 \land \phi_3 \land \phi_4 \]
\section{Determinazione interpretazioni}
\subsection{$ M \models \phi $}

Sia $D = \{\alpha, \beta \}$ il dominio di $M$:
\begin{itemize}[label=-]
    \item $ M(CaneDaCaccia) = \{ \beta \} $
    \item $ M(Gatto) = \{ \} $
    \item $ M(Persona) = \{ \alpha \} $
    \item $ M(HaTopiInCasa) = \{\} $
    \item $ M(HaIlSonnoLeggero) = \{\} $
    \item $ M(AbbaiaDiNotte) = \{ \beta \} $
    \item $ M(Possiede) = \{(\alpha, \beta)\} $
    \item $ Aldo \rightarrow \alpha $
\end{itemize}
\subsection{$ I \models \phi $}

Sia $D = \{\alpha, \beta, \gamma, \delta \}$ il dominio di $I$:
\begin{itemize}[label=-]
    \item $ I(CaneDaCaccia) = \{ \beta \} $
    \item $ I(Gatto) = \{ \delta \} $
    \item $ I(Persona) = \{ \alpha, \gamma \} $
    \item $ I(HaTopiInCasa) = \{ \gamma \} $
    \item $ I(HaIlSonnoLeggero) = \{ \gamma \} $
    \item $ I(AbbaiaDiNotte) = \{ \delta \} $
    \item $ I(Possiede) = \{(\alpha, \beta), (\gamma, \delta)\} $
    \item $ Aldo \rightarrow \alpha $
\end{itemize}
$I$ non può essere modello di $\phi$, poiché $\gamma$ possiede $\delta$, un gatto
che abbaia, pur avendo il sonno leggero e i topi in casa

\end{document}
